\section*{Chapter \ref{ch:g11:labour-market} --- The Labour Market}
% figdata: figdata/grade-11/10-labour-market/out/ (edpays, uneduc); Brenland numbers invented.

\begin{solution}{exo:g11:labour-market:1}
$3 \times$ the \omterm{def:g11:production-and-costs:marginal-product}{marginal product}: 120, 102, 84, 66 and 48 dalens an hour.
\end{solution}

\begin{solution}{exo:g11:labour-market:2}
At 70: three workers (the third adds 84, the fourth only 66). At 50: four (the fifth adds 48).
\end{solution}

\begin{solution}{exo:g11:labour-market:3}
$1578/966 \approx 1.63$; unemployment is $4.3 - 2.8 = 1.5$ points lower with a bachelor's degree.
\end{solution}

\begin{solution}{exo:g11:labour-market:4}
The job at sea: it is harder, more dangerous and far from home, so it must pay more to attract workers of the same skill. The difference is a \omterm{def:g11:labour-market:compensating-differential}{compensating differential}.
\end{solution}

\begin{solution}{exo:g11:labour-market:5}
She accepts the job at 16 and refuses the one at 12; her \omterm{def:g11:labour-market:labour-supply}{reservation wage} is 14.
\end{solution}

\begin{solution}{exo:g11:labour-market:6}
$\mathit{MPL} = 30/\sqrt{L}$, so $\mathit{VMPL} = 60/\sqrt{L}$. At 30: $\sqrt{L} = 2$, $L = 4$. At 15: $\sqrt{L} = 4$, $L = 16$.
\end{solution}

\begin{solution}{exo:g11:labour-market:7}
$\mathit{VMPL} = 90/\sqrt{L} = 30$ gives $L = 9$, up from 4. Each worker's output is now worth more, so it pays to hire workers whose output was worth less than their
wage before.
\end{solution}

\begin{solution}{exo:g11:labour-market:8}
Employers pay $w + 3$, so $2 + 0.02L + 3 = 26 - 0.04L$, $L = 350$. Workers receive $2 + 0.02 \times 350 = 9$; employers pay 12. Against the untaxed wage of 10, workers bear
1 dalen of the tax and employers 2, although the law charges it all to employers.
\end{solution}

\begin{solution}{exo:g11:labour-market:9}
$(1 + r)^4 = 1.63$ gives $1 + r = 1.63^{1/4} \approx 1.13$: about 13\,\% a year, inside the twins study's range of 12 to 16\,\%. The two may differ because graduates differ from
high-school leavers in more than schooling (ability, family), because the medians are of different people, and because a degree may be worth more than its four years.
\end{solution}

\begin{solution}{exo:g11:labour-market:10}
The \omterm{def:g11:marginal-reasoning:marginal-benefit}{marginal cost} of labour is $4 + 0.1L$; $4 + 0.1L = 28 - 0.1L$ gives $L = 120$, paid $4 + 0.05 \times 120 = 10$. Competitive: $4 + 0.05L = 28 - 0.1L$ gives $L = 160$ and
$w = 12$.
\end{solution}

\begin{solution}{exo:g11:labour-market:11}
Highest among the least educated: the United States, Britain, Germany, France, Spain and South Africa. Highest among the most educated: India (13.5\,\% against 2.3\,\%), and in
Mexico the least educated have the lowest rate. Where there is no unemployment insurance and most jobs are informal, the poor cannot afford to stay without work and take any
work there is, while graduates can wait for a formal job that matches their schooling.
\end{solution}

\begin{solution}{exo:g11:labour-market:12}
Studies in which a \omterm{def:g11:labour-market:minimum-wage}{minimum wage} rose and employment did not fall: the New Jersey fast-food restaurants compared with those of Pennsylvania, and the 138 state increases.
One side answers that the competitive model is wrong for such markets, where employers have power over wages. The other side answers that the study covers one industry and
one rise, that Pennsylvania may not show what New Jersey would have done without the rise, and that the most credible studies, in its judgement, find job losses.
\end{solution}

\begin{solution}{exo:g11:labour-market:13}
At 11: below the competitive wage, the owner hires all who will work at 11, $(11 - 4)/0.05 = 140$. At 12: 160. At 14: the owner hires where $\mathit{VMPL} = 14$, $L = 140$.
For a \omterm{def:g11:labour-market:minimum-wage}{minimum wage} $m$ between the \omterm{def:g11:labour-market:monopsony}{monopsony} wage and the competitive wage $w_c$, the owner hires the supply at $m$, which rises with $m$; above $w_c$, it hires the demand at $m$, which falls
with $m$. Employment therefore rises up to $w_c$ and falls after it, and is highest at $m = w_c$.
\end{solution}

\begin{solution}{exo:g11:labour-market:14}
Abler people may both study longer and earn more; children of richer or better-connected families may both study longer and find better jobs. Identical twins share genes
and family, so comparing twins removes both. It does not deal with whatever made one twin study longer than the other, such as motivation or health, if that also raises pay.
\end{solution}

\begin{solution}{exo:g11:labour-market:15}
Too high: the 6 unexplained points also contain every relevant difference not measured, such as interruptions of work, the features of the jobs taken, and their
\omterm{def:g11:labour-market:compensating-differential}{compensating differentials}. Too low: discrimination can act through the explained part, by keeping one group out of better-paid occupations or longer hours.
\end{solution}

\begin{solution}{pb:g11:labour-market:1}
\textbf{1.} $26 - 0.04L = 2 + 0.02L$: $L = 400$, $w = 10$.

\textbf{2.} Employment is demand at $m$, $(26 - m)/0.04$; the number who want to work is $(m - 2)/0.02$. At 11: 375 and 450; at 13: 325 and 550; at 15: 275 and 650.

\textbf{3.} Employment changes by $-25$, $-75$ and $-125$; unemployment is 75, 225 and 375.

\textbf{4.} $13 \times 325 = 4225$ dalens an hour, against $10 \times 400 = 4000$: the wage bill rises.

\textbf{5.} Between 325 and 400 workers, the value of the \omterm{def:g11:production-and-costs:marginal-product}{marginal product} exceeds the supply price; at 325 the gap is $13 - 8.5 = 4.5$: the loss is $\tfrac12 \times 75 \times
4.5 = 168.75$ dalens an hour.

\textbf{6.} Bill $L(2 + 0.02L) = 2L + 0.02L^2$; \omterm{def:g11:marginal-reasoning:marginal-benefit}{marginal cost} of labour $2 + 0.04L$.

\textbf{7.} $2 + 0.04L = 26 - 0.04L$: $L = 300$, paid $2 + 0.02 \times 300 = 8$.

\textbf{8.} The workers: 100 fewer have jobs, and those who do are paid 8 instead of 10. Each produces $26 - 0.04 \times 300 = 14$ and is paid 8, 6 below.

\textbf{9.} Under the \omterm{def:g11:labour-market:minimum-wage}{minimum wage}, the owner pays $m$ to everyone up to the number who will work at $m$, so one more worker costs exactly $m$, without raising the others' wage.

\textbf{10.} At 11 the owner hires up to $\mathit{VMPL} = 11$, 375, which is within the 450 who will work at 11: 375. At 13: 325. At 15: 275.

\textbf{11.} Competitive: $-25$, $-75$, $-125$. \omterm{def:g11:labour-market:monopsony}{Monopsony} (from 300): $+75$, $+25$, $-25$.

\textbf{12.} At the competitive wage, 10: employment is 400.

\textbf{13.} Employment under \omterm{def:g11:labour-market:monopsony}{monopsony} is $(26 - m)/0.04$ above 10; it falls below 300 when $26 - 0.04 \times 300 < m$, that is above 14.

\textbf{14.} 125 fewer under the competitive model, 25 fewer under \omterm{def:g11:labour-market:monopsony}{monopsony}.

\textbf{15.} At 15 both predict fewer jobs; at 11 and 13 they disagree on the sign.

\textbf{16.} \omterm{def:g11:labour-market:monopsony}{Monopsony}: the competitive model predicts fewer jobs for any binding \omterm{def:g11:labour-market:minimum-wage}{minimum wage}.

\textbf{17.} Whether a café that raises its wage a little draws many more applicants (a flat supply to the café: little power) or few (a steep one: power); how many employers
the town has; whether workers quit quickly when paid less.

\textbf{18.} In a small town a few employers face the whole local supply of workers and have power over wages; in a large city workers can move between many employers.

\textbf{19.} A small rise may stay below the point where even a monopsonist cuts jobs; a large one passes it under both models.

\textbf{20.} At 11, 13 and 15 dalens the competitive model predicts $-25$, $-75$ and $-125$ jobs; the \omterm{def:g11:labour-market:monopsony}{monopsony} model predicts $+75$, $+25$ and $-25$. Its prediction turns
negative for \omterm{def:g11:labour-market:minimum-wage}{minimum wages} above 14 dalens.
\end{solution}
